\section{总结与延伸}

\subsection{核心知识脉络}

本节课建立了一条从生物视觉到人工视觉的完整认知链：

\begin{center}
\begin{tikzpicture}[node distance=0.8cm, every node/.style={draw, rounded corners, minimum width=3.5cm, minimum height=0.7cm, font=\small}]
\node (bio) {生物视觉的进化};
\node (cv) [below=of bio] {计算机视觉的诞生};
\node (nn) [below=of cv] {神经网络与反向传播};
\node (data) [below=of nn] {大数据时代（ImageNet）};
\node (dl) [below=of data] {深度学习革命（AlexNet）};
\node (modern) [below=of dl] {现代 AI 爆发};
\draw[->, thick] (bio) -- (cv);
\draw[->, thick] (cv) -- (nn);
\draw[->, thick] (nn) -- (data);
\draw[->, thick] (data) -- (dl);
\draw[->, thick] (dl) -- (modern);
\node[draw=none, right=0.8cm of bio, text width=5.5cm, font=\scriptsize] {寒武纪大爆发、Hubel \& Wiesel、层级视觉通路};
\node[draw=none, right=0.8cm of cv, text width=5.5cm, font=\scriptsize] {Roberts (1963)、MIT Summer Vision Project (1966)、David Marr};
\node[draw=none, right=0.8cm of nn, text width=5.5cm, font=\scriptsize] {Perceptron (1958)、Neocognitron (1980)、Backprop (1986)、LeNet (1998)};
\node[draw=none, right=0.8cm of data, text width=5.5cm, font=\scriptsize] {Caltech-101、PASCAL VOC、ImageNet (2009)};
\node[draw=none, right=0.8cm of dl, text width=5.5cm, font=\scriptsize] {AlexNet (2012)、VGG、GoogLeNet、ResNet};
\node[draw=none, right=0.8cm of modern, text width=5.5cm, font=\scriptsize] {Transformer、扩散模型、VLM、具身智能};
\end{tikzpicture}
\end{center}

\subsection{关键 Takeaways}

\begin{importantbox}{五条核心原则}
\begin{enumerate}
    \item \textbf{视觉是智能的基石}：5.4亿年的进化告诉我们，视觉感知是驱动智能发展的核心力量。理解视觉就是理解智能。
    \item \textbf{数据与算法同等重要}：深度学习的突破不仅仅依赖于更好的算法（反向传播），还依赖于大规模的标注数据（ImageNet）。数据是深度学习的一等公民。
    \item \textbf{生物学启发贯穿始终}：从 Hubel \& Wiesel 的感受野到卷积神经网络的设计，神经科学的发现一直是算法创新的重要灵感来源。
    \item \textbf{三大力量缺一不可}：计算力、算法和数据的同时就绪，才引爆了深度学习革命。单独依靠任何一个维度都不够。
    \item \textbf{技术进步必须伴随伦理反思}：AI 的强大能力带来了偏见、公平性和社会影响等严肃议题，技术开发者有责任正视这些挑战。
\end{enumerate}
\end{importantbox}

\subsection{展望}

CS231N 接下来的课程将从图像分类和线性分类器开始，系统地构建深度学习在计算机视觉中的完整知识体系。2025年的课程还新增了大规模分布式训练、扩散模型、视觉语言模型等反映当前研究前沿的内容。正如 Fei-Fei Li 所说，我们正处于``AI 全球变暖''时期——一个激动人心且看不到放缓迹象的技术加速时代。

\end{document}
